\section{Conclusion}
\label{sec:conclusion}

We asked whether the human-\versus-AI direction that detectors read from the residual stream also controls how AI-like a model's writing sounds. In \qwen it does.
Subtracting \dai lowers supervised detector scores (\desklib 0.996 $\to$ 0.84) and raises the zero-shot \binoculars score. Equal-norm random directions, formality, verbosity, the Assistant axis, and a ``write like a human'' prompt fail to match this at equal coherence, and the effect holds among perfectly fluent outputs and after projecting out six confounds.
The direction is better described as the style signature of chat tuning than as ``machine text''. It is distinct from the Assistant persona, and it already exists, with the same orientation, in the base model, which it steers in both directions.
It is also a weak lever. Fluent steered text remains far from human scores, in-distribution clamping barely moves detectors, and an LLM judge finds prompting equally effective.

{\bf Takeaway:} ``sounding like AI'' has a real, causal linear component, but no single direction captures it.

\para{Future work.} We plan to steer with a low-rank subspace of the human--AI difference instead of one vector, and to apply steering only after the first sentence. We also plan to replicate on Gemma-2 with SAE features and on Llama-3.1-8B, to run human evaluation of fluent steered outputs, and to test whether the detectors' remaining signal lives in formatting by combining \dai with markdown removal.
